\chapter{Social Engineering}
\label{ch:1.1}
\chaptertopic{Topic 1.1}

\chapterquote{The human factor is truly security's weakest link.}{Kevin Mitnick, \emph{The Art of Deception} (2002), ch.\ 1}

\section{Motivation}

The fundamental question of this chapter: how does an attacker defeat a defense that is, by any
technical measure, unbreakable?

Consider a lock that would take a thief a thousand years to pick, guarding a door an employee is
holding open for a stranger carrying a box of donuts. That is not a hypothetical. It is how most
organizational security actually fails. Money goes into cryptography, firewalls, and access
control -- a perimeter that a competent social engineer bypasses simply by talking to a person.
Social engineering exploits no flaw in software. It exploits the decision-making shortcuts every
human brain runs under time pressure, stress, or social discomfort, shortcuts that exist because
constant, careful deliberation is not sustainable at scale.

Attackers did not discover a flaw. They found a feature, and learned to trigger it on command.

This chapter treats that exploitation as three separable mechanisms. A defense effective against
one does nothing against another.
\begin{enumerate}
\item \textbf{Pressure} -- the psychological tactics that degrade a target's judgment before any
request is made.
\item \textbf{Extraction and delivery} -- elicitation, malicious links, and malicious files: the
mechanisms that convert degraded judgment into a compromised account or system.
\item \textbf{Channel} -- email, SMS, social media: each carrying a distinct default trust
assumption that attackers borrow rather than construct.
\end{enumerate}
By the end of the chapter, a student should be able to classify any social-engineering attempt
along all three dimensions, and identify the single verification habit that defeats every
combination of them.

\begin{apalignment}{Unit 1: Introduction to Security}{Topic 1.1 -- Understanding Social Engineering}
\begin{apobjectives}
\apobjective{1.1.A}{Analyze Risk}{Recognize the common indicators that a message or interaction is a social engineering attempt.}
\apobjective{1.1.B}{Analyze Risk}{Explain how social engineering tactics persuade a victim to carry out the action an adversary wants.}
\apobjective{1.1.C}{Analyze Risk}{Describe the consequences a victim may face after a successful social engineering attack.}
\end{apobjectives}
\end{apalignment}

\section{The Four Psychological Levers}

Call this the \textbf{pressure model}: four tactics account for nearly every case of
psychologically-driven compromise on record, because human decision-making under pressure varies
remarkably little across culture, age, or education.

\begin{enumerate}
\item \textbf{Urgency} manufactures an artificial deadline. Phrases such as ``your account will be
closed in 24 hours'' are not primarily about the stated deadline; they are designed to deny the
target time to think. The deliberate, effortful component of cognition that would ordinarily
verify a sender address or inspect a link requires time to engage -- remove the time, and
verification does not occur.
\item \textbf{Authority} borrows legitimacy the sender has not earned. Compliance with figures who
appear to hold institutional power -- a supervisor, a government agency, an IT administrator -- is
a norm instilled from childhood. An attacker who convincingly impersonates such a figure inherits
the target's learned deference without performing any of the work that deference is meant to
reward.
\item \textbf{Scarcity} exploits loss aversion. Claims such as ``only 3 spots remaining'' or
``offer expires tonight'' activate the same fear of missing out that shapes ordinary consumer
behavior, redirected toward prompting action before the target applies routine skepticism.
\item \textbf{Intimidation} substitutes fear for urgency's time constraint. Threats of legal
action, account suspension, or being reported create a desire to eliminate the threat as quickly
as possible, with compliance presented as the fastest available resolution.
\end{enumerate}

\begin{definition}{Social engineering}
The manipulation of people, rather than machines or code, into taking an action that benefits an
attacker: revealing a credential, clicking a link, transferring money, or granting physical or
system access.

\emph{Informally:} an attack against the operator, not the system.
\end{definition}

What do all four levers have in common? None of them touches a keyboard. Four tactics, zero
exploits, one target: the part of the brain that trusts too fast. The defense that beats all four
will not be technical either -- the final section of this chapter names it.

The four levers are also rarely deployed in isolation. A phishing email claiming to originate from
a bank (authority), warning of suspicious activity with a one-hour deadline (urgency and
intimidation together), and offering a single ``secure verification link,'' stacks three levers at
once. Stacking compounds pressure rather than simply adding to it.

\begin{example}[Why the tactic matters more than the template]
Attackers constantly change surface details -- the impersonated brand, the wording, the delivery
channel -- but the underlying psychological lever rarely changes. A defender who has memorized
``watch out for emails claiming to be from the IRS'' is defenseless against a never-before-seen
message using the same authority lever through a different brand. A defender who instead asks what
a message is trying to make them feel, and why it wants them to act quickly, recognizes the
manipulation regardless of its surface dressing -- including AI-generated attacks that have never
been seen before (Chapter~\ref{ch:1.4}).
\end{example}

\section{Elicitation: Extracting Information Through Conversation}

Pressure fails against a target who does not panic. What is the attacker's next avenue?

Call this the \textbf{extraction model}: it inverts the pressure model's logic. Pressure pushes a
target toward a mistake; extraction invites disclosure instead. Pressure alone never compromises a
system. The second phase is extraction -- getting the target to hand over information -- or
delivery -- getting the target to open something dangerous. Elicitation covers extraction.

\begin{definition}{Elicitation}
The covert extraction of sensitive information through seemingly casual conversation, structured
so the target does not realize they are being questioned for a purpose.

\emph{Informally:} interrogation disguised as small talk.
\end{definition}

Every elicitation gambit obeys a single design constraint: a direct question invites suspicion, so
a functioning gambit never asks directly. Four indirect approaches recur:

\begin{itemize}
\item \textbf{Flattery and appeal to expertise} invites the target to demonstrate knowledge at
length, revealing architectural or process details in the course of showing off.
\item \textbf{Deliberate incorrect statements} exploit the near-universal impulse to correct an
error -- the target volunteers the accurate information unprompted.
\item \textbf{Mutual disclosure} establishes an implied social debt: the attacker shares something,
true or fabricated, first, and the target feels obliged to reciprocate.
\item \textbf{Third-party reference} invokes a trusted name (``your colleague told me...'') to
frame compliance as an action that is already normal and already endorsed.
\end{itemize}

Each individual disclosure feels harmless in isolation, which is precisely what makes elicitation
difficult to detect. An attacker assembles a complete picture across several unremarkable
conversations, none of which registers, to the target, as a security event. The next two sections
examine what an attacker does once pressure or elicitation has produced an opening: deliver a
payload directly, through a link or a file.

\section{Malicious Links and Files}

Where elicitation extracts information through conversation, malicious links and files deliver a
payload directly. Both reduce to the same underlying principle: the attacker controls what the
target perceives, and constructs that perception to imply safety.

\subsection{Reading a URL like an attacker does}

A trusted brand name appearing anywhere within a URL establishes nothing about ownership. The only
component of a URL that determines ownership is the registered domain immediately preceding the
first slash, and the reliable method for locating it is to read the string \emph{backward from that
slash}, not left to right.

Five technical patterns recur throughout malicious links. \textbf{Typosquatting} substitutes a
visually similar character (\texttt{paypa1.com}). \textbf{Homograph attacks} substitute look-alike
characters from another script entirely. \textbf{Userinfo disguise} places a trusted name before an
\texttt{@}, where it is merely a discarded username rather than part of the host. \textbf{IP-literal
URLs} send a victim to a raw address instead of a hostname. \textbf{Mismatched link text} is the odd
one out: the visible anchor text names one destination while the underlying \texttt{href} points
somewhere else entirely, and because that mismatch lives in the surrounding HTML rather than in the
URL string itself, it cannot be caught by inspecting the string alone.

\begin{lstlisting}[language=Python, caption={A single function, tested against four of the five patterns above}]
from urllib.parse import urlparse

def registered_domain(url: str) -> str:
    """Registered domain a URL points to.
    The only part an attacker cannot fake with
    a trusted prefix. Cost: O(len(url))."""
    host = urlparse(url).netloc
    return host.split("@")[-1].split(":")[0]

samples = [
    "http://microsoft.com.evil.net/reset",  # subdomain label
    "https://paypa1.com/login",  # typosquat
    "https://accounts.google.com@evil.tk/verify",  # userinfo
    "http://192.168.94.12/verify-account",  # IP literal
]
for url in samples:
    print(f"{url:46s} -> {registered_domain(url)}")
# Every row: the printed domain does not match the brand read first.
\end{lstlisting}

In every row, the value returned does not match the brand a reader perceives at first glance -- in
the third row, ``accounts.google.com'' is not part of the domain at all; it is a discarded
username, and the request goes to \texttt{evil.tk}. That gap, between the domain a reader
recognizes and the domain that actually receives the request, is the entire mechanism of this
class of attack.

\begin{figure}[h]
\centering
\begin{tikzpicture}[
  seg/.style={thick, minimum height=0.9cm, align=center, font=\ttfamily\small},
  lbl/.style={font=\scriptsize\sffamily, align=center, text width=2.4cm}
]
\node[seg, draw=rulegray, fill=black!4, text=black!55, minimum width=3.1cm]
  (sub) at (0,0) {microsoft.com};
\node[seg, draw=kickstarternavy, fill=kickstarternavy!12, text=kickstarternavy, font=\ttfamily\small\bfseries, minimum width=1.6cm]
  (dom) at (2.35,0) {evil};
\node[seg, draw=kickstarternavy, fill=kickstarternavy!12, text=kickstarternavy, font=\ttfamily\small\bfseries, minimum width=1.0cm]
  (tld) at (3.75,0) {.net};
\node[seg, draw=kickstarterochre, fill=kickstarterochre!10, text=black!70, minimum width=1.9cm]
  (path) at (5.35,0) {/reset};
\draw[rulegray] (-1.55,0.45) -- (6.3,0.45);
\node[lbl, black!55] at (0,1.35) {attacker-chosen label\\ means nothing};
\node[lbl, kickstarternavy] at (3.05,1.6) {\textbf{registered domain}\\ the only part that\\ determines ownership};
\draw[-{Latex[length=2mm]}, kickstarternavy, thick] (3.05,1.05) -- (3.05,0.5);
\node[lbl, kickstarterochre!85!black] at (5.9,1.35) {path -- never\\ part of ownership};
\draw[-{Latex[length=2mm]}, kickstarterochre, thick] (5.6,1.0) -- (5.45,0.5);
\draw[{Latex[length=2mm]}-, kickstarterforest] (6.3,-0.9) -- (-1.55,-0.9);
\node[font=\scriptsize\sffamily, kickstarterforest] at (2.35,-1.25) {read backward from the first slash};
\end{tikzpicture}
\caption{Anatomy of \texttt{http://microsoft.com.evil.net/reset}: everything before the registered
domain is an attacker-chosen label that carries no ownership meaning at all.}
\label{fig:1.1-url-anatomy}
\end{figure}

\subsection{Recognizing dangerous files}

File-based attacks exploit default operating-system behavior rather than a defect in it: the
operating system does exactly what it was configured to do. Windows hides known file extensions by
default, so an attacker can name a file \texttt{invoice.pdf.exe}; the file browser displays only
\texttt{invoice.pdf}, and a user expecting a document double-clicks what is actually an executable.
Directly executable formats (\texttt{.exe}, \texttt{.scr}, \texttt{.ps1}, \texttt{.vbs},
\texttt{.js}) run arbitrary code the instant they are opened. Office documents with macros
(\texttt{.docm}, \texttt{.xlsm}) are disabled by default in modern software, which is precisely why
attackers add explicit social pressure -- ``you must enable content to view this document'' --
asking the victim to personally override a protection built for their own benefit.

\begin{misconception}
\textbf{``HTTPS means the site is safe.''} HTTPS (Hypertext Transfer Protocol Secure) certifies
that the connection is encrypted, not that the destination is trustworthy. Free, automated
certificate issuance means attackers obtain valid TLS (Transport Layer Security) certificates for
malicious domains routinely; the padlock icon says nothing about who controls the site behind it.
\end{misconception}

\section{Case Studies in Social Engineering}

The tactics described above are abstractions until observed against real organizations. The four
incidents below span credential phishing, a real-time multi-factor bypass, a purely text-based
fraud with no technical exploit at all, and a ten-minute phone call that disabled a casino's
operations. In each case, identify the specific human decision that, had it gone differently,
would have interrupted the entire chain -- that decision point is always present, and locating it
is the analytical skill this section develops.

\begin{example}[The 2016 DNC Breach]
In mid-2016 the threat group tracked as APT28 ran a spear-phishing campaign against the Democratic
National Committee and Hillary Clinton campaign staff, including campaign chairman John Podesta.
The emails impersonated Google security alerts, warning of an unauthorized login attempt and
offering a ``Change Password'' button that led to a look-alike domain built to match Google's login
page pixel for pixel. Staff who entered their real passwords on the spoofed page handed them
directly to the attackers, who used the credentials (protected by no meaningful multi-factor
authentication at the time) to exfiltrate over 50,000 internal emails later published online. A
now-famous detail of the incident: an aide asked to check the email's legitimacy typed ``This is a
legitimate email'' when he meant to type the opposite -- a single word, undoing whatever scrutiny
had been applied, and Podesta clicked the link. One human error, on top of one technical deception,
was sufficient.\srccite{1.1}{1}
\end{example}

\begin{example}[The July 2020 Twitter Vishing Incident]
A group of attackers used open-source reconnaissance to map Twitter's internal support staff, then
placed voice-phishing (vishing) calls impersonating Twitter IT, claiming a VPN configuration update
required employees to log in through a look-alike portal. The portal proxied the login in real
time: as an employee typed their password, the attacker's fake page relayed it instantly to the
real Twitter login, which then prompted for a one-time multi-factor code. The victim, still
believing they were completing a routine VPN update, typed that code into the fake portal too, and
the attackers relayed it to the real site within its short validity window. With that access, the
attackers reached an internal administrative tool, took control of 130 high-profile accounts, and
used them to post a cryptocurrency scam that collected \$118{,}000 in Bitcoin within
minutes.\srccite{1.1}{2} This is the
practical limit of time-based one-time-password multi-factor authentication: it defeats a stolen,
static credential, but a real-time relay attack simply forwards each factor before it expires.
Hardware security keys built on WebAuthn/FIDO2 (Web Authentication / Fast IDentity Online 2) close
this gap because the cryptographic challenge is bound to the exact domain requesting it -- a
look-alike portal cannot obtain a valid signature no matter how convincing it looks to the human on
the other end.
\end{example}

\begin{example}[Business Email Compromise: Toyota Boshoku, 2019]
Not every social-engineering loss involves a stolen credential. In August 2019 a European
subsidiary of Toyota Boshoku Corporation wired approximately \$37 million to an attacker-controlled
account after a finance employee received an email -- apparently from a trusted business partner --
requesting an urgent change to the bank details for a pending wire transfer. The message used
accurate, specific knowledge of a real transaction already in progress, almost certainly harvested
from an earlier, unrelated mailbox compromise, and combined it with the standard authority-plus-urgency
pairing. No malware, no phishing link, and no forged domain were involved: the message may even
have originated from the partner's genuine, already-compromised mailbox, which is exactly why
domain-authentication controls like DMARC (Domain-based Message Authentication, Reporting, and
Conformance) do nothing to stop it. Business email compromise of this kind is consistently the
highest-dollar-value social-engineering category precisely because it requires no exploit -- only a
believable pretext and a deadline. The only reliable control is procedural rather than technical:
any change to payment instructions is verified through a phone number looked up independently,
never one supplied in the email itself.\srccite{1.1}{3}
\end{example}

\begin{example}[Help-Desk Vishing: MGM Resorts, 2023]
Some of the most damaging social-engineering attacks target not the eventual victim but the
people whose job is to help them. In September 2023, a group known as Scattered Spider built a
profile of an MGM Resorts employee from information posted on LinkedIn -- name, job title, and
enough biographical detail to sound credible. A caller then contacted MGM's IT help desk,
impersonated that employee, and claimed to be locked out of their account. No credential was
stolen and no phishing link was clicked; the entire attack was a ten-minute phone call. The help
desk, following a process built for speed rather than verification, reset the employee's
credentials and granted the caller access without confirming identity through any channel the
attacker did not already control. That single reset was enough: within days the attackers had
escalated to broad administrative access across MGM's systems, and on September 11, 2023 they
deployed ALPHV/BlackCat ransomware, forcing casino floors, hotel room keys, and reservation
systems offline for over a week. MGM's own regulatory filing reported an approximately \$100
million hit to Adjusted Property EBITDAR for the third quarter of 2023, plus under \$10 million in
one-time incident-response costs. Company leadership separately signaled a further \$30--40 million
in forward-looking cybersecurity capital investment for the following year.\srccite{1.1}{4} The
vulnerability exploited was not technical -- it was a help-desk
procedure that treated a confident, well-informed voice on the phone as sufficient proof of
identity. This is why organizations that face high-value vishing risk increasingly require a
callback to a number on file, or an in-person or video identity check, before any credential
reset -- the same ``verify through a path the attacker does not control'' principle that closes
phishing also closes vishing, just applied to a help desk instead of an individual employee.
\end{example}

\section{Delivery Channels: Email, SMS, and Social Media}

The psychological levers and delivery mechanisms described above are channel-agnostic: they operate
identically regardless of how a message reaches its target. Each delivery channel, however, carries
its own default trust assumption, and attackers select channels deliberately for the assumption
that channel provides at no cost.

\textbf{Email} remains dominant because the underlying protocol, SMTP (Simple Mail Transfer
Protocol), was never designed with sender authentication built in -- any server can claim to send
mail as any address. Modern countermeasures (SPF, DKIM, and DMARC -- Sender Policy Framework,
DomainKeys Identified Mail, and Domain-based Message Authentication, Reporting, and Conformance,
respectively) close this gap only where a domain has adopted and correctly configured them.
\textbf{Phishing} is the mass variant, sent to many recipients with no specific target in mind.
\textbf{Spear phishing} is researched and tailored to one person or organization. \textbf{Whaling}
is spear phishing aimed specifically at executives, whose authority can approve large transfers or
grant broad access.

\textbf{SMS phishing (smishing)} reduces to the same four levers and the same delivery mechanism as
email phishing, aimed through a channel that carries weaker default scrutiny: the personal, trusted
feel of a text message, a screen too small for comfortable URL inspection, and no sender domain to
check at all. \textbf{Social media attacks} exploit a third trust model entirely: a message that
appears to come from someone already in the target's network, or a brand they already follow, faces
a lower bar of scrutiny than an unsolicited email from a stranger.

\begin{table}[h]
\centering
\begin{tabular}{@{}ll@{}}
\toprule
Channel & Default trust assumption exploited \\
\midrule
Email & ``This looks like it's from an organization I deal with.'' \\
SMS & ``Texts are personal and direct; this is from someone who has my number.'' \\
Social media & ``This is from someone in my network, or a brand I follow.'' \\
\bottomrule
\end{tabular}
\caption{Each delivery channel borrows a different kind of default trust.}
\end{table}

Sophisticated campaigns increasingly span channels: a ``pre-warning'' text (``expect an important
email from us shortly'') primes the target before the actual phishing email arrives, manufacturing
social proof across two channels that, individually, might each have raised suspicion.

\begin{misconception}
\textbf{``I can always tell a phishing email by its grammar.''} Spear phishing crafted from real
research, and increasingly AI-generated phishing (see Chapter~\ref{ch:1.4}), can be
grammatically flawless and contextually precise. Grammar quality is not a reliable defense.
\end{misconception}

\section*{The One Defense That Generalizes}

Every technique in this chapter -- pressure tactics, elicitation, malicious links and files,
channel-specific trust -- reduces to the same structural countermeasure: \textbf{verify through a
path the message itself did not provide.} Do not call the number supplied in the text. Do not use
the link in the email to ``log in and check.'' Locate the organization's real contact information
independently, through a channel the attacker does not control, and confirm the request there.

This one practice defeats nearly every variant in this chapter, without requiring the defender to
recognize which specific scam is in play. It works because it denies the attacker any influence
over the verification path -- the one structural feature every version of this attack shares.

Elicitation resists this fix only partially. There is no path to verify, because the danger is not
a false claim arriving through a channel -- it is an unremarkable conversation that never announces
itself as an attack. Its defense is different in kind: treat any specific operational detail as
sensitive by default, regardless of who is asking or how casually, since by the time a question
feels worth verifying, the information has usually already been given away.

Chapter \ref{ch:1.4} extends this analysis to attacks generated by AI systems rather than human
operators. The levers stay identical. Only the scale changes, by orders of magnitude.

\section*{Review Questions}
\begin{reviewquestions}
\item A message claims to be from your company's CEO, states that a wire transfer must be sent
within the hour to close an urgent deal, and asks you not to loop in the finance team. Identify
every psychological lever in play and explain what each one is designed to prevent you from
doing.
\item Explain why reading a URL ``right to left from the first slash'' is a more reliable
technique than scanning the URL for a trusted brand name.
\item A colleague says, ``I always check that emails don't have spelling mistakes before trusting
them.'' Explain why this heuristic is unreliable, and propose two indicators that generalize
better.
\item Compare the default trust assumption exploited by SMS phishing to the one exploited by
email phishing. Why might a user who has been trained to scrutinize email links still fall for a
smishing message?
\item Describe the mutual-disclosure elicitation gambit and explain why the information given up
during it often does not feel, to the target, like a security-relevant disclosure.
\item In the MGM Resorts vishing case, no malware was used and no link was clicked. Explain why
this attack still counts as social engineering, and identify the specific procedural failure that
allowed a ten-minute phone call to lead to a full ransomware deployment.
\item A help desk wants to stop vishing attacks like the one that hit MGM without making
legitimate employees wait longer for routine password resets. Propose a verification step that
closes the gap the attacker exploited, and explain why it does not simply move the same
vulnerability somewhere else.
\end{reviewquestions}

\begin{answerkey}
\item A wire-transfer deadline framed as coming from the CEO stacks three levers at once.
Authority: the message borrows the CEO's standing, so the request is not questioned the way a
peer's would be. Urgency: the one-hour deadline denies the target time to engage careful,
deliberate judgment. Isolation: the instruction not to loop in finance removes the one colleague
positioned to notice the request is irregular. Each lever is designed to prevent the same
underlying action: independent verification through a channel the message itself does not control.
\answerlabel{1.1.B}
\item Reading right-to-left from the first slash finds the registered domain -- the only part of a
URL an attacker cannot spoof by adding a trusted name elsewhere in the string. Scanning left to
right for a familiar brand name is unreliable instead, because that name can be placed anywhere: as
a subdomain label, a path segment, or plain text. None of those placements establish ownership of
the domain that actually receives the request.
\answerlabel{1.1.A}
\item Grammar quality is not a reliable signal, because spear-phishing built from real
reconnaissance, and AI-generated phishing in particular, can be flawless and contextually precise.
Two indicators generalize better. First: whether the message creates pressure to act before
verifying -- urgency, intimidation, or scarcity. Second: whether the request or link asks the
recipient to authenticate or disclose information through a path the message itself supplies,
rather than one the recipient already trusts independently.
\answerlabel{1.1.A}
\item Email phishing exploits the assumption that a message looking like it is from a known
organization can be trusted. SMS phishing exploits a different assumption: that a text is personal
and direct, because it required knowing the recipient's phone number. A user trained to hover over
links and check sender domains in email has often received no equivalent training for SMS. The
small screen discourages URL inspection, and there is no visible sender domain to check at all.
The trust model is different, so the defense has to be different too, even though the underlying
levers are the same.
\answerlabel{1.1.A}
\item In mutual disclosure, the attacker shares something first -- true or fabricated -- to create
an implied social debt, and the target reciprocates by sharing something comparable. The exchange
feels like ordinary conversational give-and-take rather than an interrogation. Because of that, the
target does not experience the disclosure as answering a question at all, and each fragment
volunteered seems too small and too informal to register as a security-relevant event.
\answerlabel{1.1.C}
\item It still counts as social engineering because the attacker's goal was to manipulate a human
decision -- the help-desk agent's choice to reset credentials -- rather than to exploit a technical
flaw in software. The pretext (a confident, well-informed impersonation built from LinkedIn
research) and the psychological lever (a locked-out employee who needs urgent help) are the same
structural elements as phishing, just delivered by voice. The procedural failure was narrower and
more specific: the help desk accepted claimed identity, supported only by information the attacker
could plausibly have gathered from public sources, as sufficient grounds to reset credentials, with
no verification through a channel the caller did not control.
\answerlabel{1.1.B}
\item A callback to the phone number already on file for that employee -- not a number the caller
provides -- verifies identity through a channel the attacker does not control. That is exactly
what was missing. This does not relocate the vulnerability: the attacker would need to have
already compromised that specific, independently-looked-up number to intercept the callback, and
that is a categorically harder task than researching a name and title on LinkedIn. A step that
instead asked the caller to ``confirm'' information as proof of identity would fail for a simpler
reason -- attackers routinely gather exactly that kind of information beforehand.
\answerlabel{1.1.B}
\end{answerkey}

\section*{Sources}
\begin{enumerate}[label=\lbrack\arabic*\rbrack, leftmargin=2.2em]
\item \srcitem{1.1}{1}\emph{United States v.\ Viktor Borisovich Netyksho, et al.}, No.\ 1:18-cr-00215 (D.D.C.\ filed
July 13, 2018) -- the Department of Justice indictment detailing the GRU's spear-phishing operation
against the DNC and John Podesta.
\item \srcitem{1.1}{2}Help Net Security, ``High-Profile Twitter Accounts Hijacked to Push Bitcoin Scam,'' July 16,
2020; Bitcoin figure per public blockchain analysis of the attackers' wallet.
\item \srcitem{1.1}{3}Lee Mathews, ``Toyota Parts Supplier Hit By \$37 Million Email Scam,'' \emph{Forbes},
September 6, 2019.
\item \srcitem{1.1}{4}MGM Resorts International, Form 8-K, filed with the U.S.\ Securities and Exchange Commission,
October 5, 2023.
\end{enumerate}
